\subsection{Objekt KL47}
\label{subsec:ObjektKL47}


\small

\paragraph{1. Objektsteckbrief}

\par Bei \gls{kl47} handelt es sich um ein \gls{mfh} in der Karl-Liebknecht-Straße 47, 07546 Gera. Es handelt sich um ein ursprünglich im Jahr $XXX$ errichtetes \gls{mfh} das derzeit als , das in den Jahren 1992 bis 1994 kernsaniert und modernisiert wurde und seither instand gehalten wird. Die \gls{mfh} verfügt über $XXX$ Wohneinheiten und $XXX$ Gewerbeeinheiten von denen sich $XXX$ im Eigentum der \gls{stz} befindet. Die vermietbare Wohnfläche liegt bei $XXX m^2$, die vermietbare Gewerbefläche liegt bei $XXX m^2$.

\begin{table} [htbp!]
\begin{tabular*}{\textwidth}{@{\extracolsep{\fill}}>{\raggedright\arraybackslash}p{4.5cm}>{\raggedright\arraybackslash}p{10.3cm}}
\toprule
\textbf{Merkmal} & \textbf{Angabe} \\
\midrule
Adresse & Karl-Liebknecht-Straße 47, 07546 Gera \\
Objektart & \gls{go} \\
Baujahr / Modernisierung & 1994 \\
Vermietbare Fläche & $XXX m^2$ Wohnen \\
Einheiten & $X$ Wohneinheiten \\
		& $X$ Gewerbeeinheiten \\
Vermietungsstand & 100\% per \today \\
Eigentümer & \gls{stz} \\
\bottomrule
\end{tabular*}
\caption{Objektsteckbrief KL47}
\label{tab:objektsteckbriefKL47}
\end{table}

\paragraph{2. Ertragslage und Stabilität (p.a.)}
\begin{table} [htbp!]
\begin{tabular*}{\textwidth}{@{\extracolsep{\fill}}>{\raggedright\arraybackslash}p{8.0cm}>{\raggedleft\arraybackslash}p{2.2cm}>{\raggedleft\arraybackslash}p{2.2cm}>{\raggedleft\arraybackslash}p{2.2cm}}
\toprule
\textbf{Kennzahl} & \textbf{2023} & \textbf{2024} & \textbf{2025} \\
\midrule
Nettokaltmiete 	& XX 		&  XX 	&  XX \\
Nicht umlagefähige Bewirtschaftungskosten & XX & XX & - \\
NOI (Net Operating Income) & 13.100 & 28.882 & - \\
Leerstandsquote & 0\% & 0\% & 0\% \\
\bottomrule
\end{tabular*}
\caption{Ertragslage H25a}
\label{tab:ertragslageH25a}
\end{table}


\paragraph{3. Bankrelevante Kennzahlen (stabilisierter Zustand)}
\begin{table} [htbp!]
\begin{tabular*}{\textwidth}{@{\extracolsep{\fill}}>{\raggedright\arraybackslash}p{6.1cm}>{\raggedleft\arraybackslash}p{2.8cm}>{\raggedright\arraybackslash}p{6.3cm}}
\toprule
\textbf{Kennzahl} & \textbf{Wert} & \textbf{Einordnung} \\
\midrule
Marktwert / interner Ansatz & [EUR] & Bewertungsstichtag: [Datum] \\
Beleihungswert (konservativ) & [EUR] & Sicherheitsabschlag: [\%] \\
Wunschdarlehen auf Objektbasis & [EUR] & Anteil am Gesamtobligo: [\%] \\
LTV (Darlehen / Marktwert) & [\%] & Bankziel: [z.B. \textless= 70\%] \\
Debt Yield (NOI / Darlehen) & [\%] & Bankziel: [z.B. \textgreater= 8\%] \\
DSCR (CFADS / Schuldendienst) & [x,xx] & Bankziel: [z.B. \textgreater= 1,20] \\
\bottomrule
\end{tabular*}
\caption{Bankrelevante Kennzahlen H25a}
\label{tab:bankkennzahlenH25a}
\end{table}

\paragraph{4. Risiko- und Sicherheitenprofil}
\begin{itemize}
	\item \textbf{Mietausfall-/Leerstandsrisiko:} Ein leerstabiles Objekt mit 100\% Vermietungsquote, abgesichert durch langfristige Mietverträge mit bonitätsstarken Mietern.
	\item \textbf{Instandhaltungsrisiko:} Derzeit keine größeren Instandhaltungsmaßnahmen erforderlich, Rücklagenbildung für zukünftige Maßnahmen gemäß Instandhaltungsplan.
	\item \textbf{Marktrisiko:} Lage in Sollngriesbach mit stabilem Immobilienmart, jedoch noch knapp 30 jährige Ölheizung im Objekt, was bei zukünftigen Vermietungen zu berücksichtigen ist.\endnote{Eine Umrüstung auf eine nachhaltige Heizlösung wird mittelfristig angestrebt. Hierzu sind bereits erste Gespräche mit Fachfirmen geführt worden.}
	\item \textbf{Rechtliche Sicherheiten:} Es besteht ein  Nießbrauchrecht der vorherigen Eigentümer
\end{itemize}

\paragraph{5. Kurzfazit und Ausblick}
Das Bestandsobjekt Hopfengasse 25a zeigt einen sehr stabilen Zustand. Es existiert eine tragfähige Cashflow-Basis im Objekt, dass derzeit für die \{gls} aufgrund des bestehenden Nießbrauchrechts keine Möglichkeit zur Bedienung von Zins und Tilgung möglicher Kredite.

\vspace{0.5em}
\textit{Hinweis: Detailrechnungen (Tilgungsplan, Sensitivitäten, Mieterliste) sind im Anhang bzw. in den Tabellenblättern hinterlegt.}

\normalsize
